\documentclass[11pt,a4paper]{article}
\usepackage[margin=21mm,headheight=15pt]{geometry}
\usepackage[T1]{fontenc}
\usepackage{lmodern,amsmath,amssymb,booktabs,tabularx,enumitem,xcolor,fancyhdr}
\usepackage[hidelinks]{hyperref}
\definecolor{accent}{RGB}{0,114,178}
\setlength{\parindent}{0pt}
\setlength{\parskip}{6pt}
\setlist{nosep,leftmargin=*}
\pagestyle{fancy}
\fancyhf{}
\fancyhead[L]{\small Econometrics 25117 | Instructor guide}
\fancyhead[R]{\small 2026--2027}
\fancyfoot[L]{\small Private teaching notes}
\fancyfoot[R]{\thepage}
\newcommand{\heading}[1]{\section*{\color{accent}#1}}
\newcommand{\cue}[2]{\par\noindent\textbf{#1}\quad #2\par}
\setlength{\emergencystretch}{1em}
\newcommand{\E}{\mathbb{E}}
\newcommand{\Var}{\operatorname{Var}}
\newcommand{\Cov}{\operatorname{Cov}}
\begin{document}
\begin{center}{\Large\bfseries Threats to identification}\end{center}
\textbf{Slide route:} Lecture8v2.tex: external/internal validity, measurement error, missing data, selection, simultaneity, SEs and prediction. Chapter 9.
\heading{Session 11 | Diagnose before choosing a remedy}
\textbf{Outcomes:} Separate internal from external validity; explain classical attenuation, selection, and simultaneity; match threats to plausible responses.
\heading{100-minute route}
\begin{tabularx}{\textwidth}{@{}rX@{}}\toprule
0--10 & Post-midterm retrieval: interpretation versus identification.\\
10--25 & External and internal validity; discuss slide examples.\\
25--45 & OVB, wrong functional form, measurement error.\\
45--55 & Pause and attenuation calculation.\\
55--75 & Missing data, selection, simultaneity.\\
75--90 & Inconsistent SEs, validity checklist, prediction.\\
90--100 & Exit diagnosis and bridge to IV.\\\bottomrule
\end{tabularx}
\heading{Worked example | Measurement error}
\cue{Use with slides}{Lecture8v2 slides 12--15: Errors-in-variables and Classical Measurement Error. Work the attenuation calculation on slide 14 or 15, after students have seen why the measured regressor contains noise.}
Suppose the true relationship is $Y=\beta_0+2X^*+u$, but we observe $X=X^*+v$. Under classical regressor measurement error, $v$ is mean zero and uncorrelated with $X^*$ and $u$. If $\Var(X^*)=9$ and $\Var(v)=3$:
\[
\operatorname{plim}\hat\beta_{\rm OLS}
=2\frac{9}{9+3}=1.5.
\]
\cue{Ask}{Does a larger sample restore the coefficient to 2? No. It concentrates the estimator around the attenuated target. Better measurement or a valid instrument addresses a different problem from increasing precision.}
\cue{Qualification}{Attenuation toward zero is a result for this simple classical setup; do not generalize its direction to arbitrary measurement error or every coefficient in multiple regression. Classical noise in the outcome alone need not bias slopes, though it raises uncertainty.}

\newpage
\heading{Cases for discussion | Threat, mechanism, response}
\cue{Selection}{Wages are observed only among employed people. If employment depends on unobserved wage determinants, the observed sample need not represent the wage-offer relationship for everyone. More complete data or a justified selection design is needed; merely adding a robust SE is insufficient.}
\cue{Simultaneity}{Price and quantity are jointly determined by demand and supply. Their observed covariance need not trace a demand curve. A supply shifter excluded from demand is a possible instrument, provided its exclusion and relevance can be defended.}
\cue{External validity}{An internally valid tutoring experiment at one university need not identify the same effect for all schools. Ask what differs: students, implementation, context, or treatment intensity.}
\cue{Dependence}{Repeated outcomes for one district can be correlated. Heteroskedasticity-robust SEs alone are insufficient; a dependence-aware variance estimator may be needed.}
\heading{Pair activity | Classify three claims}
\begin{enumerate}
\item A treatment improves outcomes in one city, so it must do so everywhere.
\item Self-reported education contains substantial recall error.
\item A high-$R^2$ demand regression uses equilibrium prices.
\end{enumerate}
\textbf{Answers:} External validity; potential errors-in-variables; simultaneity despite fit. For each, ask what extra data or design evidence would make the claim more credible.
\cue{Missing data caution}{Dropping incomplete rows is not automatically harmless. Whether complete-case analysis is valid depends on why data are missing and the estimand. Avoid treating all missingness as one problem with one remedy.}
\cue{Prediction distinction}{An endogenous variable can be predictive in a stable environment, yet policy interventions can change its relationship to the outcome. Evaluate predictions out of sample; evaluate causal claims through design and assumptions.}
\cue{If short / preparation}{Use one external-validity illustration and three internal threats in depth. Read Chapter 12 and come prepared to propose an instrument and criticize its exclusion restriction.}
\textbf{Source:} Lecture8v2.tex; numerical and discussion examples are illustrative.

\end{document}
